%==========================================================================
% CONCLUSION (unnumbered)
%==========================================================================
\unnumberedchapter{General Conclusion and Perspectives}\label{ch:conclusion}

\section*{Summary of the Work}

This thesis set out to design a telemonitoring system for patients with hypertension and type~2 diabetes in Cameroon, and to evaluate its predictive component without the biases of previous work. The literature showed that telemonitoring improves blood pressure and glycated haemoglobin when a clinician acts on the data, that the only validated protocol for patients with both conditions triages single readings with fixed thresholds, and that the predictive results of the field rest on labels that leak into the features and on validation in the training population (\cref{ch:review}). The risk-tiering problem was therefore formalised as a forecast of the next reading, labelled by official thresholds, combined with the guideline rule so that a model can raise an alert but never lower it (\cref{ch:framework}). A pipeline of data collection and treatment and a benchmark protocol were built to select the model rather than presuppose it, and the system was built around them (\cref{ch:methodology}).

\section*{Principal Findings}

The data work produced four results. First, the Douala charts were taken with electronic devices for the pressures, while the pulse was read off the chart curve; the readings themselves showed this, since the devices are not named. Second, the device tolerance alone can change the level of a quarter to more than half of the readings, which bounds what any model can achieve against a threshold label. Third, the sourced labelling rule changes the level of 58 of the 366 readings compared with the earlier unsourced rule, chiefly because low pressures now count and the glucose bands follow the official thresholds. Fourth, the 32 local patients are far below the size needed to develop a model, and the development data differ from them in prevalence and in the frequency of critical hyperglycaemia, which predicts a loss of calibration that a local intercept update should partly recover.

The benchmark answered the predictive questions with a mixed result. On blood pressure, learned models forecast the level of the next reading better than the guideline rule applied to the last reading: five of eight families, and the selection procedure as a whole, improve on persistence significantly, in most hospitals and in every subgroup, but by about 3\% of the ranked probability score, and the deployed CatBoost relies mostly on the last systolic pressure. On glucose, no model beats persistence significantly, on a task that the demonstration data cannot fully support. Deep temporal models did not outperform tree ensembles, although GRU-D was the procedure's consistent choice for glucose. Applied to the Douala patients, every model lost accuracy and under-predicted the alert; an intercept update corrected the offset but not the slopes, the value forecasts arrived nearly calibrated without it, and no model improved the alert decisions over the rule. The measures of the evolution of the patient's state carry little information about the next level, and the one retained is shown to the clinician as a description, not used as an alert. The lead times of 24 and 48~hours were not tested.

These results justify the design choice made before any of them was known: the guideline rule is the floor, and a model can only add to it.

The system answers the first research question: a patient application and a clinician dashboard that implement the three-level home protocol offline, a record that cannot be altered, a queue that orders the clinician's attention, and server functions for escalation, reminders and verdicts, verified by 219 automated tests, among them the parity of the triage computed on the phone with the Python reference.

\section*{Future Directions}

\begin{itemize}
  \item \textbf{Full development data.} Rerunning the benchmark on the full eICU database, available under a data use agreement, removes the sample-size constraint of the glycaemic task.
  \item \textbf{Clinical input.} The relabelling of 100 to 150 readings by the partner physician, and the costs of errors that fix the weights of the selection cost and the tier cut-points, replace the provisional choices.
  \item \textbf{Home data.} A prospective pilot in Douala would collect the home readings for which the system is designed, revalidate the model on them, and measure adherence, usability and the clinician's workload; the clinician verdicts it records become labels that do not derive from the thresholds.
  \item \textbf{Acquisition.} Readings from Bluetooth devices would remove transcription errors, and a low-cost sensor node (ESP32 with pulse oximetry and temperature) could add covariates; neither was built in this work.
  \item \textbf{Deployment.} Hosting the risk service and the worker, an SMS fallback for patients without data, and a database region chosen under Law No.~2024/017 are required before any use with patients.
\end{itemize}
